% =====================================================================
\chapter{Polinomios, ecuaciones y sistemas}
\label{tema:3}
% =====================================================================

Los polinomios son las funciones más sencillas después de las lineales y aparecen por todas partes: como aproximaciones de Taylor, como funciones de pérdida cuadráticas y como polinomio característico de una matriz. Este tema repasa su aritmética (división, raíces, factorización), la resolución de ecuaciones e inecuaciones y los sistemas lineales pequeños. Sigue los capítulos~2, 5 y~7 de OpenStax \emph{College Algebra}~\cite{openstax}, con las demostraciones que allí se omiten.

\section{Polinomios y división euclídea}

\begin{definicion}[Polinomio]
Un \emph{polinomio} con coeficientes reales es una expresión $P(x)=a_nx^n+\dots+a_1x+a_0$ con $a_i\in\R$. Si $a_n\neq0$, $n$ es el \emph{grado} de $P$, $\deg P$, y $a_n$ su \emph{coeficiente principal}. Un número $r$ es una \emph{raíz} de $P$ si $P(r)=0$.
\end{definicion}

Si $P$ y $Q$ son no nulos, $\deg(PQ)=\deg P+\deg Q$, porque el coeficiente principal del producto es el producto de los coeficientes principales, que no es cero (\cref{prop:cuerpo}(iv)).

\begin{teorema}[División euclídea]\label{teo:division}
Sean $P$ y $D$ polinomios con $D\neq0$. Existen polinomios únicos $Q$ (cociente) y $R$ (resto) tales que
\[
P=DQ+R\qquad\text{y}\qquad R=0\ \text{o}\ \deg R<\deg D.
\]
\end{teorema}

\begin{teorema}[Teoremas del resto y del factor]\label{teo:resto}
El resto de dividir $P$ entre $x-a$ es el número $P(a)$. En consecuencia, $a$ es raíz de $P$ si y solo si $x-a$ divide a $P$.
\end{teorema}

La \emph{regla de Ruffini} es una forma abreviada de dividir entre $x-a$: se bajan los coeficientes uno a uno, multiplicando el último resultado por $a$ y sumándolo al coeficiente siguiente.

\addtocounter{definicion}{1}%

\begin{corolario}\label{cor:raices}
Un polinomio no nulo de grado $n$ tiene a lo sumo $n$ raíces reales distintas.
\end{corolario}

\begin{teorema}[Raíces racionales]\label{teo:raices-racionales}
Sea $P(x)=a_nx^n+\dots+a_0$ con coeficientes enteros, $a_n\neq0$ y $a_0\neq0$. Si $p/q$, con $p,q$ enteros primos entre sí, es raíz de $P$, entonces $p$ divide a $a_0$ y $q$ divide a $a_n$.
\end{teorema}

\section{La ecuación de segundo grado y las fórmulas de Vieta}

\begin{teorema}[Ecuación de segundo grado]\label{teo:cuadratica}
Sean $a\neq0$ y $\Delta=b^2-4ac$ el \emph{discriminante}. La ecuación $ax^2+bx+c=0$ tiene dos soluciones reales distintas si $\Delta>0$, una (doble) si $\Delta=0$ y ninguna si $\Delta<0$. Cuando $\Delta\ge0$,
\[
x=\frac{-b\pm\sqrt{\Delta}}{2a}.
\]
\end{teorema}

\begin{proposicion}[Fórmulas de Vieta]\label{prop:vieta}
Si $x_1,x_2$ son las raíces (reales o repetidas) de $ax^2+bx+c$, entonces
\[
ax^2+bx+c=a(x-x_1)(x-x_2),\qquad x_1+x_2=-\frac ba,\qquad x_1x_2=\frac ca.
\]
Más en general, si $P(x)=a_n\prod_{i=1}^n(x-x_i)$, entonces $\sum_ix_i=-a_{n-1}/a_n$ y $\prod_ix_i=(-1)^na_0/a_n$.
\end{proposicion}

\addtocounter{definicion}{1}%

\section{Otras ecuaciones e inecuaciones}

Las ecuaciones \emph{bicuadradas} $ax^4+bx^2+c=0$ se reducen a segundo grado con $t=x^2$. En las ecuaciones \emph{con radicales} se aísla la raíz y se eleva al cuadrado; como elevar al cuadrado no es una operación reversible ($a^2=b^2$ no implica $a=b$), pueden aparecer soluciones \emph{extrañas} y es obligatorio comprobar cada candidata en la ecuación original. En las ecuaciones \emph{racionales} se multiplica por el denominador común y se descartan los valores que lo anulan.

Para resolver una inecuación $P(x)/Q(x)\ge0$ se factorizan numerador y denominador, se marcan en la recta los puntos donde se anula algún factor y se estudia el signo de cada factor en cada intervalo (\emph{tabla de signos}); los ceros del denominador nunca forman parte de la solución.

\addtocounter{definicion}{1}%

\section{Sistemas lineales pequeños}

Un sistema de $m$ ecuaciones lineales con $n$ incógnitas es \emph{compatible determinado} si tiene una única solución, \emph{compatible indeterminado} si tiene infinitas e \emph{incompatible} si no tiene ninguna. Los métodos de sustitución, igualación y reducción de bachillerato son casos particulares de la eliminación gaussiana del módulo~3.

\begin{teorema}[Regla de Cramer para $2\times2$]\label{teo:cramer}
El sistema $ax+by=e$, $cx+dy=f$ tiene solución única si y solo si $\delta=ad-bc\neq0$, y entonces
\[
x=\frac{ed-bf}{\delta},\qquad y=\frac{af-ec}{\delta}.
\]
\end{teorema}

\addtocounter{definicion}{1}%

\addtocounter{definicion}{1}%

\addtocounter{definicion}{1}%

